\section{Evaluation protocol}
\label{sec:protocol}

The protocol is the paper's main methodological contribution, so we state the
rules it enforces and why each exists.

\subsection{Splitting}

Splits are by song, never by clip. A song yields four clips at \SI{5}{\second}; if clips
of one song fall on both sides, the model can recognise the song rather than the
generator. Splits are stratified by the four-way taxonomy so that rare
sub-classes are represented everywhere.

We note a property of SONICS's own published splits that appears not to have
been reported: its \texttt{id} field links a real song to the synthetic tracks
generated from it, and 4{,}349 source identifiers---12.9\% of rows---appear in
more than one official split. A real recording can therefore sit in the test set
while a track generated from its lyrics and style sits in validation. We report
this rather than repairing it, since re-splitting forfeits comparability with
published figures.

\subsection{Removing shortcuts before measuring}

Before any detector is trained we ask whether the classes are separable without
reference to generation at all, by fitting a linear model on eight global
descriptors of the recording---loudness, crest factor, spectral rolloff,
centroid, zero-crossing rate and high-frequency energy share. None of these
describe how audio was produced; a large AUROC means a 94-million-parameter
transformer will find the same shortcut immediately.

The first such test returned \textbf{0.925 AUROC}. Commercial releases are
mastered loud and heavily compressed while generated tracks are not, so loudness
and crest factor alone carried the classification. Peak normalisation reduced
this to 0.840 and root-mean-square normalisation to 0.821; crest factor is
scale-invariant and survives both. Randomising dynamic range---applied to both
classes and to evaluation as well as training, since the aim is to remove a
confound from the task rather than to augment one split---brought the residual
to \textbf{0.714}, and to \textbf{0.737} on the final 910-song corpus. We report
detector performance only against that baseline.

A second diagnostic compares the spectral bandwidth of the two classes. Real
audio fetched from streaming sources at one bitrate and generated audio
distributed at another are trivially separable on codec artefacts. We match the
bitrate distribution of the real class to the generated set and verify that the
median spectral-rolloff ratio falls below a fixed threshold before reporting any
figure.

\subsection{Three evaluation axes}

\paragraph{In-distribution.}
Held-out songs, same generators, same corpus. A sanity number.

\paragraph{Cross-generator, shared corpus.}
SONICS contains both Suno and Udio. Training on one and testing on the other
isolates the generator: corpus, pipeline, processing and real class are all held
fixed. Real songs are divided disjointly between the halves, and both directions
are run because transfer is not symmetric in general.

\paragraph{Cross-corpus.}
FakeMusicCaps pairs five unseen text-to-music generators with real music from
MusicCaps. Here the generator \emph{and} the corpus of real music change
together. We band-limit the real clips to the generated set's sample rate and
verify bandwidth before reporting.

\subsection{Reporting rules}

Every cross-generator result is reported beside the $F_1$ that a classifier
labelling all inputs synthetic would obtain on the same class balance. This is
not a formality: our own cross-corpus $F_1$ of 0.750 matches that baseline to
four decimal places. AUROC and false-positive rate are reported alongside $F_1$
throughout, because $F_1$ alone cannot distinguish a working detector from a
degenerate one.
